\chapter{Public API and MCP Tool Catalog}
\label{app:api}

\section{Python API}
\label{sec:api}

The orchestrator package exposes three symbols that make up the
platform's Python API:

\begin{lstlisting}[caption={\file{orchestrator/\_\_init\_\_.py}.},
label={lst:api}]
from orchestrator.loop import stream_chat
from orchestrator.pool import MCPClientPool, register_default_servers

__all__ = ["stream_chat", "MCPClientPool", "register_default_servers"]
\end{lstlisting}

A minimal programmatic session:

\begin{lstlisting}[caption={Minimal script using the Python API.},
label={lst:api-usage}]
import asyncio
from orchestrator import MCPClientPool, register_default_servers, stream_chat

async def main():
    pool = MCPClientPool()
    register_default_servers(pool)
    await pool.start()

    async for evt in stream_chat(
        query="Run EDA on iris.csv",
        active_file="iris.csv",
        history=[],
        pool=pool,
    ):
        print(evt)

    await pool.shutdown()

asyncio.run(main())
\end{lstlisting}

Direct tool invocation is also possible without the LLM in the
loop:

\begin{lstlisting}[caption={Direct MCP tool call (bypasses the LLM).},
label={lst:direct}]
result = await pool.call_tool(
    "ingest_dataset", {"file_path": "iris.csv"})
print(result.structured_content)
\end{lstlisting}

\section{MCP Tool Catalog}
\label{sec:tool-catalog}

Table~\ref{tab:tool-catalog} lists every tool with its owning
service, required arguments, and return summary. The full JSON
schemas are auto-generated by FastMCP from the Python type hints
and docstrings and are visible inside the running app on the
\emph{Tools} tab.

\begin{table}[H]
\centering
\caption{Complete MCP tool catalog.}
\label{tab:tool-catalog}
\small
\begin{tabularx}{\linewidth}{@{}llXX@{}}
\toprule
\textbf{Service} & \textbf{Tool} & \textbf{Required args} & \textbf{Returns (summary)} \\
\midrule
mcp-data & \tool{list\_uploaded\_files} & --- & \code{count}, \code{files[]} \\
mcp-data & \tool{ingest\_dataset} & \code{file\_path} & shape, dtypes, head, missing, duplicates \\
mcp-data & \tool{validate\_schema\_with\_pandera} & \code{file\_path} & \code{valid}, inferred rules, failure cases \\
mcp-eda & \tool{run\_full\_eda} & \code{file\_path} & stats, missing summary, 3 plot artefacts \\
mcp-eda & \tool{render\_correlation\_matrix} & \code{file\_path} & method, columns, plot artefact, strong pairs \\
mcp-modeling & \tool{run\_parallel\_bake\_off} & \code{file\_path}, \code{target\_column} & leaderboard, champion, model \& xtrain artefact IDs \\
mcp-modeling & \tool{trigger\_hyperparameter\_sweep} & \code{file\_path}, \code{target\_column}, \code{model\_family} & best value, best params, trial count \\
mcp-explain & \tool{calculate\_shap\_values} & \code{model\_id}, \code{x\_train\_sample\_id} & plot artefact, global importance, explainer kind \\
mcp-explain & \tool{generate\_feature\_importance\_plot} & \code{model\_id} & plot artefact, top-$k$ features \\
mcp-export & \tool{generate\_jupyter\_notebook} & \code{file\_path}, \code{target\_column}, \code{problem\_type} & notebook artefact, cell count \\
mcp-export & \tool{compile\_pdf\_report} & \code{project\_title}, \code{file\_path}, \code{target\_column}, \code{leaderboard}, \code{champion\_model} & PDF artefact \\
\bottomrule
\end{tabularx}
\end{table}

\section{MCP Prompts}
\label{sec:prompts}

The platform ships with one prompt template:

\begin{itemize}[leftmargin=1.4em]
  \item \tool{eda\_deep\_dive(file\_path, target\_column="")} ---
        instructs the LLM to call \tool{run\_full\_eda},
        \tool{render\_correlation\_matrix}, and summarise strong
        correlations, skewed features, and next-step recommendations.
\end{itemize}
